\section{S1\quad DFT Protocol}

\subsection*{Input}
All reference calculations use ORCA 6.0.1. Each structure is computed in two runs. The first is a UKS single point with stability analysis:

\begin{verbatim}
! UKS wB97M-V def2-TZVPD def2/J RIJCOSX TightSCF DEFGRID3
%scf
  Thresh 1e-12
  TCut   1e-13
  MaxIter 300
  STABPerform true
  STABRestartUHFifUnstable true
end
* xyzfile 0 1 ts.xyz
\end{verbatim}

The second computes energy and forces (\texttt{EnGrad}) and reads the converged orbitals of the first run (\texttt{MORead}). A gradient run from a fresh guess could converge back to the closed-shell solution. Its forces would then belong to a different solution than the one the spin class was determined on.

\subsection*{ORCA version}
OMol25 used ORCA 6.0.0\cite{omol25}. In that version, the stability analysis reports spurious negative eigenvalues for $\omega$B97M-V; in a test on 17 structures, 9 runs aborted. ORCA 6.0.1 corrects this. OMol25 itself was not affected, because it did not use the stability analysis.

\subsection*{Agreement with the OMol25 labels}
We recomputed 44 Transition1x transition states with this protocol, and also with ORCA 5.0.4 (Table~S1). Both versions give the same spin class as OMol25 for all 44 structures (27 closed-shell, 17 broken-symmetry). The energy deviates from the OMol25 label by a nearly constant 0.07\,meV, at most 0.1\,meV. The largest force component deviates by at most 0.013\,eV/\AA{}.

\begin{table}[h]
  \caption{Deviation of our protocol from the OMol25 labels and between ORCA versions, at 44 Transition1x transition states. Force: difference of the largest force component.}
  \label{tab:S1}
  \begin{tabular}{lcccc}
    \hline
    & \multicolumn{2}{c}{Energy (meV)} & \multicolumn{2}{c}{Force (eV/\AA)} \\
    Comparison & median & max & median & max \\
    \hline
    ORCA 6.0.1 vs OMol25      & 0.072  & 0.098 & 0.0004 & 0.013 \\
    ORCA 6.0.1 vs ORCA 5.0.4  & 0.0001 & 0.028 & 0.0006 & 0.007 \\
    \hline
  \end{tabular}
\end{table}

% Numbers: results/orca6_gate.csv (44 rows, dE6omol_mev, dF6omol, dE65_mev, dF65).
% TODO 5.0.4 row only if the thesis/5.0.4 is mentioned
% TODO "9 of 17" is from D16 (docs/paper_decisions.md); needs a results file for the checker
% TODO the 38 lower solutions (table) go to S2 next to the label check

\section{S2\quad Reaction Selection}
Candidates were chosen with $N_\mathrm{FOD}$ and the OMol25 $\langle S^2\rangle$, which is available for every transition state.
% TODO: depth bins (1-10, 10-43, 43-300, >300 meV: 30/34/33/33), twin rule (within 0.05, same heavy-atom
% count where possible, no reuse), seed 0, candidate pools (s2pos 1,271, cshigh 656, twins 364, base 200)

\section{S3\quad Transition-State Search}
The search follows Marks et al.\cite{marks2026fsm} and uses ML-FSM 1.0.1\cite{marks2026fsm} and Sella 2.6.0\cite{hermes2022sella} with ASE.

\paragraph{Endpoints.} The Transition1x reactant and product are relaxed on each model with BFGS to a largest atomic force below 0.05\,eV/\AA{} (at most 2,000 steps). Overall translation and rotation between the two are removed before the search.

\paragraph{Freezing string.} The string is grown from both ends with 18 nodes, interpolated in redundant internal coordinates. Each new node is optimized perpendicular to the path with L-BFGS-B in Cartesian coordinates (at most 2 optimization steps with 3 line-search steps each, step length at most 0.05\,\AA). The node of highest model energy is the transition-state guess.

\paragraph{Saddle-point refinement.} Sella refines the guess to a first-order saddle point in internal coordinates, until the largest atomic force is below 0.05\,eV/\AA{}, with at most 1,000 steps. Runs that do not converge are not used for the error analysis.

\begin{table}[h]
  \caption{Settings of the transition-state search.}
  \begin{tabular}{ll}
    \hline
    Endpoint relaxation & BFGS, 0.05\,eV/\AA, $\leq$ 2,000 steps \\
    Freezing string     & 18 nodes, redundant internal coordinates \\
    Node optimization   & L-BFGS-B, 2 steps, 3 line-search steps, max.\ step 0.05\,\AA \\
    Saddle refinement   & Sella, order 1, internal coordinates, 0.05\,eV/\AA, $\leq$ 1,000 steps \\
    \hline
  \end{tabular}
\end{table}
% TODO check the meaning of the CartesianOptimizer arguments ('L-BFGS-B', 2, 3, 0.05) in ML-FSM;
% docstring of fsm_search.py says "at most 3 line-search and 2 optimisation steps per node, step limit 0.05 A".
% TODO eSEN endpoints come from the NEB runs (omol_neb.py, same BFGS settings) -- check steps limit there.
% Source: ~/opt/mlfsm_src (thegomeslab, commit f1f53a0), ~/venvs/sella (sella 2.6.0)
